% Informe de arbitraje y auditoría de replicación.
% Compilar dentro de esta carpeta:  tectonic referee_report.tex
\documentclass[11pt,a4paper]{article}
\usepackage[margin=2.4cm]{geometry}
\usepackage{fontspec}
\setmainfont{texgyretermes}[Extension=.otf,UprightFont=*-regular,BoldFont=*-bold,
  ItalicFont=*-italic,BoldItalicFont=*-bolditalic]
\setmonofont{lmmono10-regular.otf}[Scale=0.92]
\usepackage{unicode-math}
\setmathfont{texgyretermes-math.otf}
\usepackage[spanish,es-nodecimaldot,es-noshorthands,es-tabla,es-nolists,es-nosectiondot]{babel}
\usepackage{booktabs,longtable,array,tabularx,xcolor,enumitem,caption}
\usepackage{xurl}
\usepackage[colorlinks=true,linkcolor=black,urlcolor=blue!45!black]{hyperref}

\setlength{\parindent}{0pt}
\setlength{\parskip}{5pt plus 1pt}
\setlist{itemsep=2pt,topsep=3pt,parsep=0pt,leftmargin=1.5em}
\captionsetup{font=small,labelfont=bf,justification=raggedright,singlelinecheck=false,skip=4pt}
\renewcommand{\arraystretch}{1.12}
\emergencystretch=2em

\newcommand{\coment}[2]{\medskip\noindent\textbf{#1.\enspace #2}\par\nopagebreak}
\newcommand{\pide}{\par\noindent\textbf{Qué se pide.}\enspace}
\newcommand{\cod}[1]{\texttt{#1}}
\newcolumntype{L}[1]{>{\raggedright\arraybackslash}p{#1}}
\newcolumntype{Y}{>{\raggedright\arraybackslash}X}

\begin{document}

\begin{center}
{\Large\bfseries Informe de arbitraje y auditoría de replicación}\\[6pt]
{\large \emph{Internet Shutdowns and Software Development: Evidence from Bangladesh}}\\[10pt]
\begin{tabular}{@{}ll@{}}
Autor del manuscrito: & Hector Alberto Diaz Gomez (borrador de octubre de 2026) \\
Revisión solicitada por: & Alexander Quispe (asesor) \\
Versión auditada: & commit \cod{7346f22}, \path{paper/latex/main.tex} \\
Fecha: & 7 de octubre de 2026 \\
\end{tabular}
\end{center}

\smallskip
{\small Este informe se elaboró con Claude Code ejecutando el repositorio completo. Toda cifra que no
proviene del manuscrito sale de los scripts de \path{audit_scripts/}, cuyas salidas están en
\path{audit_scripts/output/}. Los cálculos marcados como \emph{diagnóstico} son verificaciones de
arbitraje; no son especificaciones nuevas del paper.\par}

\section{Resumen}

\textbf{Replicación.} Las cifras del manuscrito se reproducen exactamente a partir de los datos
procesados que están en el repositorio: control sintético, placebos, diferencias en diferencias, RDiT y
mecanismo; 20 de 22 tablas y las 8 figuras se regeneran idénticas. El paquete todavía no es replicable
desde los datos crudos, y los outcomes secundarios tienen un error de muestra que cambia sus magnitudes
en puntos porcentuales.

\textbf{Arbitraje.} El resultado principal es creíble en signo y orden de magnitud. En su forma actual
no pasaría una revista top-5 por cuatro razones: (i) el hallazgo central es casi mecánico, porque sin
Internet no se puede publicar en GitHub; (ii) la evidencia de trabajo offline, leída en niveles, indica
que la creación de código casi se detuvo, lo contrario del énfasis del texto; (iii) hay un quiebre en
la serie de Bangladesh el 2 de julio, antes del tratamiento, que afecta al contrafactual; y (iv) la
inferencia está limitada por un grupo de 14 donantes que el autor eligió.

\textbf{Lo que está bien hecho.} Los scripts \cod{01\_reproduce\_*} detienen la ejecución si no
reproducen el valor archivado. Las tablas se generan por código. El lenguaje es prudente y el valor
$1/15$ está bien explicado. La ponderación de los placebos del mecanismo está documentada y auditada.

\textbf{Cinco prioridades, en orden.}
\begin{enumerate}
\item Corregir los outcomes secundarios (R1).
\item Rehacer el mecanismo por fecha de autoría y en niveles (M1).
\item Fijar la elegibilidad antes del período de estimación y tratar el quiebre de julio (M2).
\item Ampliar el grupo de donantes y la inferencia (M4).
\item Reencuadrar las contribuciones; llevar el análisis sectorial y el RDiT al apéndice (M6 a M8).
\end{enumerate}

\textbf{Regla de trabajo.} Ningún cambio debe adoptarse porque produzca resultados más fuertes. Varios
de los pedidos reducen magnitudes. Cuando una especificación cambie, hay que reportar la original junto
a la nueva. Los cambios de M2 y M4 son metodológicos mayores: conviene acordarlos con el asesor antes de
implementarlos.

\section{Auditoría de replicación}

\subsection{Procedimiento}

Copié el repositorio a un directorio temporal y ejecuté todo ahí, sin tocar los archivos originales.
Entorno: Python 3.12, pandas 3.0.6, numpy 2.5.3, scipy 1.18.1, pyarrow 25.0.1, matplotlib 3.11.2;
\cod{tectonic} para LaTeX. Ejecuté, en este orden, los tres scripts de
\path{analysis/synthetic_control/}, y los de \path{analysis/did/}, \path{analysis/rdit/},
\path{analysis/mechanism/}, \path{analysis/tables/} y \path{analysis/figures/}. Comparé las salidas con
lo versionado: archivos CSV y \cod{.tex} byte a byte, figuras píxel a píxel a 150 dpi, y el texto del
PDF recompilado contra \path{paper/latex/thesis.pdf}.

\subsection{Estado por componente}

\begin{table}[h]
\small
\begin{tabularx}{\textwidth}{@{}L{6.2cm}Y@{}}
\toprule
Componente & Resultado \\
\midrule
Control sintético, placebos entre países, leave-one-out & Reproduce exacto: 66 verificaciones, 0 fallos \\
Placebos en el tiempo, DiD, event study, RDiT, placebos del mecanismo & Reproducen exacto \\
SDID & Solo se recalcula con los pesos archivados; no se reestima \\
Tablas LaTeX & 20 de 22 idénticas; 2 fallan (R3) \\
Figuras & 8 de 8 idénticas; la A1 no se regenera sola (R3) \\
PDF & Compila sin referencias rotas; texto idéntico al PDF versionado \\
Panel país-día del outcome principal & Lo reconstruí exacto desde \path{data/processed/}; el código no está en el repo \\
Tablas principales del mecanismo & Las reconstruí exactas; el código no está en el repo \\
Outcomes secundarios & Reconstruí las tablas archivadas, pero tienen un error de muestra (R1) \\
Análisis sectorial & No replicable: los insumos están en Google Drive \\
Extracción de GH Archive, elegibilidad, filtro de bots, muestreo de metadata & Sin código (R6) \\
\bottomrule
\end{tabularx}
\end{table}

\subsection{Hallazgos}

\coment{R1}{Los outcomes secundarios mezclan dos muestras}
En \path{results/synthetic_control/country_day_outcomes_enriched.parquet} el numerador cuenta a los
150,017 usuarios del archivo usuario-día, y el denominador solo a los 59,158 elegibles. Las 1,905
celdas país-día coinciden exactamente con esa construcción, y no con la de solo elegibles. El numerador incluye,
por tanto, cuentas no elegibles y las excluidas como bots. La tabla \ref{tab:r1} compara lo publicado
con mi reconstrucción usando solo elegibles. Los efectos relativos casi no cambian; los puntos
porcentuales bajan entre 17\% y 37\%.

\begin{table}[h]
\small
\caption{Outcomes secundarios: manuscrito frente a reconstrucción con usuarios elegibles}
\label{tab:r1}
\begin{tabularx}{\textwidth}{@{}Yrrrrrr@{}}
\toprule
 & \multicolumn{2}{c}{Manuscrito} & \multicolumn{2}{c}{Corregido} & \multicolumn{2}{c}{Rango placebo corregido} \\
\cmidrule(lr){2-3}\cmidrule(lr){4-5}\cmidrule(l){6-7}
Outcome & ATT & Relativo & ATT & Relativo & ATT & Ratio RMSPE \\
\midrule
\multicolumn{7}{@{}l}{\emph{A. Pesos primarios fijos}} \\
Productiva (pp) & −15.179 & −86.5\% & −11.665 & −86.6\% & & \\
Shipping (pp) & −1.333 & −94.2\% & −0.979 & −93.3\% & & \\
Colaboración (pp) & −0.780 & −93.6\% & −0.611 & −93.5\% & & \\
Eventos de colaboración / 1,000 & −24.095 & −95.1\% & −19.956 & −95.5\% & & \\
Shipping output / 1,000 & −23.981 & −93.9\% & −17.913 & −93.0\% & & \\
\addlinespace
\multicolumn{7}{@{}l}{\emph{B. Pesos específicos por outcome}} \\
Productiva (pp) & −14.870 & −86.2\% & −11.720 & −86.7\% & 0.067 & 0.067 \\
Shipping (pp) & −0.964 & −92.2\% & −0.668 & −90.5\% & 0.067 & 0.067 \\
Colaboración (pp) & −0.630 & −92.1\% & −0.480 & −91.9\% & 0.067 & 0.133 \\
Eventos de colaboración / 1,000 & −18.589 & −93.7\% & −11.640 & −92.5\% & 0.267 & 0.133 \\
Shipping output / 1,000 & −18.192 & −92.1\% & −13.274 & −90.8\% & 0.067 & 0.067 \\
\bottomrule
\end{tabularx}
\par\smallskip{\footnotesize \emph{Notas:} «Corregido» usa las mismas columnas del archivo usuario-día
(\cod{productive\_output\_active}, \cod{shipping\_active}, \cod{collaboration\_active} y los conteos)
restringidas a los elegibles, con el estimador de \path{01_reproduce_sc.py}. Es mi reconstrucción; la
versión definitiva debe salir del código original. Fuente: \path{03_secondary_sc_reproduce.py}.\par}
\end{table}

El DiD por outcome tiene el mismo problema: para actividad productiva pasa de −14.669 a −11.748 pp.
Además, el manuscrito no reporta rangos placebo para los dos outcomes de conteo. Con el panel archivado
dan 0.200 y 0.200 para eventos de colaboración, y 0.133 y 0.067 para \emph{shipping output}.
\pide Reconstruir el panel de outcomes secundarios solo con elegibles; regenerar
\path{results/final/table_secondary_*.csv} y \path{table_DiD_all_outcomes.csv}; actualizar las tablas
\cod{tab:secondary}, \cod{tab:did\_secondary} y \cod{tab:weights\_specific}, la sección 6.4 y la frase
de la introducción con −86.2, −92.2 y −92.1. Reportar los rangos placebo de los outcomes de conteo.
Subir el código al repositorio.

\coment{R2}{El outcome primario y la elegibilidad no están definidos de forma reproducible}
El outcome primario es \cod{productive\_events > 0}. Esa variable no es la suma de ninguna combinación
de las columnas del archivo: incluye tipos de evento que no figuran ni en los datos ni en el texto. De
los 1,103,680 usuario-día elegibles con actividad «productiva», 104,786 (9.5\%) no tienen push, pull
request, issue, release, review ni comentario. La Tabla 2 del manuscrito define «productive activity»
de otra manera (commits, issues cerrados, pull requests cerrados o releases); con esa definición la
tasa previa de Bangladesh es 15.38\% y no 17.25\%.

La elegibilidad («at least four pre-treatment active days») cuenta días con \emph{cualquier} evento,
no días productivos. Con esa regla hay 60,729 usuarios; el archivo de elegibles tiene 59,158. Los 1,571
restantes se excluyeron con una regla de bots que no puedo reconstruir: las columnas \cod{name\_bot} y
\cod{behavior\_bot} valen cero para todos los elegibles.
\pide Listar en la sección 4.2 los tipos de evento de GH Archive que componen el outcome primario.
Escribir la regla de elegibilidad y la de bots con sus umbrales. Indicar que los días son UTC (ver M5).
Subir el código que construye \path{eligible_users_donors.parquet}.

\coment{R3}{Los scripts de tablas y figuras fallan en un clon limpio}
\begin{itemize}
\item \path{analysis/tables/make_all_tables.py} se detiene en \cod{tab\_sector\_recovery}: lee
\path{results/sector_recovery/sectorwide/sectorwide_support.csv}, que está en \cod{.gitignore}.
\cod{tab\_sector\_fixed\_repo} falla igual con \path{pre8/sector_recovery_summary_pre8.csv}.
\item \path{analysis/figures/make_all_figures.py} define \cod{fig\_sc\_intime\_placebos} pero
\cod{main()} no la llama. Llamada a mano funciona.
\item \path{analysis/sector_recovery/10_make_sector_results_final.py} falla: falta
\path{sectorwide_070/sectorwide_user_eventstudy.csv}.
\end{itemize}
\pide Versionar esos archivos pequeños o eliminar la dependencia, y llamar a la figura A1 desde
\cod{main()}. Criterio de aceptación: ambos scripts terminan sin error en un clon nuevo.

\coment{R4}{Cifras del texto que no se pueden verificar con el repositorio}
Sección 4.4: 23,866 cuentas, 168,875 repositorios, 92,979 README recuperados, 32,491 repositorios
clasificados. Sección 8.3: el 82.5\% de \emph{Arts}. Tabla \cod{tab:sector\_fixed\_repo}: las cuatro
columnas de repositorios fijos, y con ellas el rango «10.1 a 31.4» de la sección 6.5.
\pide Guardar en \path{results/sector_recovery/} un archivo pequeño con esos conteos y el resumen
\cod{pre8}.

\coment{R5}{Conflicto sobre el 5 de agosto}
La nota de la Figura 1 dice que las caídas del 17 de junio y del 5 de agosto «coincide with national
holidays». El panel codifica el 5, 6 y 7 de agosto como «Public Holiday». Pero
\path{data/processed/panel_clean.parquet}, del análisis previo del autor, codifica el 4 y 5 de agosto
como \cod{f4\_second} y el 28 de julio al 3 de agosto como \cod{f3\_socialblock}. El 5 de agosto
Bangladesh registra 4.93\% frente a 13.89\% del sintético. El manuscrito no menciona un segundo corte.
\pide Documentar con fuentes qué ocurrió entre el 4 y el 7 de agosto [CITATION NEEDED] y corregir la
nota de la Figura 1 y el párrafo de la sección 6.5 que dice que los datos «do not identify the source
of this second-week decline»: la semana +2 (31 de julio a 6 de agosto) contiene esos días.

\coment{R6}{Falta el código aguas arriba}
No están en el repositorio: las consultas a GH Archive y la construcción de
\path{data/processed/*.parquet}; la obtención de las listas de usuarios por país; la agregación a
país-día y la codificación de feriados; la estimación por outcome y sus placebos; la estimación SDID
(su error estándar no está documentado); la grilla de sensibilidad del RDiT; la extracción de
PushEvents, el muestreo diario de commits anteriores al 18 de julio (sin semilla), la consulta de
metadata y las tablas \path{table_mechanism_*.csv}; y la selección de la muestra de 20,000 cuentas
(sin semilla). Los scripts 01, 03 y 04 de esta auditoría reconstruyen tres de esos pasos y pueden
servir de base.

\coment{R7}{Empaquetado}
No hay \cod{requirements.txt} para el análisis principal. Tres archivos de \path{data/processed/}
guardan la fecha con el tipo \cod{dbdate} y exigen el paquete \cod{db-dtypes}. \path{README.md} dice
que el manuscrito principal es \path{paper/thesis.docx} y \path{main.tex} dice que es
\path{paper/thesis.txt}; ambos están desactualizados frente al LaTeX. El \path{manifest.csv} omite 33
archivos versionados.

\subsection{Cifras del texto}

Crucé más de cien cifras escritas a mano en el resumen y las secciones contra los archivos de resultados.
Coinciden, con dos observaciones menores. La sección 6.5 dice que el índice de \emph{Information} está
«between 111.5 and 126.4» en las semanas +10 a +15; el máximo es 126.35. La misma sección dice que
\emph{Real Estate} supera 100 en las semanas +5 y +6; también lo supera en la +7 (108.3).

\section{Informe de arbitraje}

\subsection{Resumen del manuscrito}

El paper estima el efecto del apagón nacional de Internet de Bangladesh (18 al 23 de julio de 2024)
sobre la actividad pública en GitHub de 10,241 cuentas asociadas al país. Con un control sintético
sobre 14 países, la proporción diaria de cuentas con actividad cae de un contrafactual de 15.28\% a
2.02\%: −13.26 pp, −86.8\%. Añade dos piezas: commits con fecha de autoría durante el apagón que se
publican después, y una descripción de la recuperación por sector económico del repositorio.

\subsection{Evaluación general}

La caída existe y es enorme; ninguna de las pruebas que hice cambia eso. El problema para una revista
top-5 es de contribución. El efecto sobre publicar en GitHub cuando no hay Internet es casi una
identidad. Las preguntas con contenido económico son otras: si la producción se perdió o se desplazó,
si hubo recuperación posterior, y si el costo persiste. El paper las plantea en las contribuciones dos
y tres, que son las de evidencia más débil. Con el rediseño que sigue puede ser un buen artículo de
revista de campo.

\subsection{Comentarios mayores}

\coment{M1}{El mecanismo: la participación es alta, el nivel es mínimo}
El 28.3\% promedia dos poblaciones. La ventana «18 al 23 de julio» en UTC empieza quince horas antes
del corte total, que Belson (2024) sitúa cerca de las 15:00 UTC del 18. La tabla \ref{tab:m1} separa
la cohorte.

\begin{table}[h]
\small
\caption{Commits de la muestra estricta con fecha de autoría del 18 al 23 de julio}
\label{tab:m1}
\begin{tabularx}{\textwidth}{@{}Yrrrr@{}}
\toprule
Fecha de autoría (UTC) & Commits & Usuarios & Publicados desde el 24 & Rezago \\
\midrule
18 de julio & 1,939 & 722 & 153 & 7.9\% \\
19 de julio & 182 & 77 & 139 & 76.4\% \\
20 de julio & 182 & 79 & 149 & 81.9\% \\
21 de julio & 179 & 69 & 133 & 74.3\% \\
22 de julio & 172 & 66 & 125 & 72.7\% \\
23 de julio & 278 & 101 & 130 & 46.8\% \\
\addlinespace
Antes de las 15:00 UTC del 18 & 1,835 & 692 & 77 & 4.2\% \\
Desde las 15:00 UTC del 18 & 1,097 & 250 & 752 & 68.6\% \\
\addlinespace
Total (manuscrito) & 2,932 & 835 & 829 & 28.3\% \\
\bottomrule
\end{tabularx}
\par\smallskip{\footnotesize Fuente: \path{04_mechanism_reproduce.py}.\par}
\end{table}

El 63\% de los commits «creados durante el apagón» se creó antes del corte, con conectividad; 1,758
de esos 1,835 se publicaron antes del 24. Entre los creados ya sin conexión, el rezago es 68.6\%.

El nivel importa más. Las seis ventanas placebo de seis días tienen en promedio 33,787 commits de la
muestra estricta, ponderados. Durante el corte total hubo 1,097: el 3.2\% de lo normal. Los 829
commits del rezago equivalen al 2.5\%. Es decir, la creación de código que termina siendo pública
cayó cerca de 97\%, y el trabajo offline recuperó una fracción muy pequeña. La frase del resumen
(«local work followed by delayed synchronization») es correcta para quienes siguieron trabajando,
pero la introducción («a substantial margin of work») y la discusión («displaced across time rather
than destroyed») no se sostienen en niveles. Advertencia: las ventanas placebo son del 1 al 11 de
julio, cuando la actividad estaba inusualmente alta (M2); aun reduciendo esa base a la mitad la
conclusión no cambia.
\pide
\begin{enumerate}
\item Definir la cohorte desde el inicio del corte total y reportar ambas definiciones.
\item Presentar la serie diaria de commits por fecha de autoría (con los pesos de muestreo antes del
18 de julio) como outcome propio: es la medida más cercana a producción que tiene el paper.
\item Reescribir resumen, introducción, sección 7, discusión y conclusión con el nivel además de la
participación.
\item Discutir la censura: los pushes se observan hasta el 10 de agosto, y 35 commits del rezago
llegan entre el 5 y el 10.
\item Explicar los 345 commits creados durante el corte total y publicados antes del 24 (64
usuarios; 246 el 23 de julio). Los publicados entre el 18 y el 22 indican cuentas con conectividad:
fuera del país o con acceso especial.
\end{enumerate}

\coment{M2}{Quiebre previo al tratamiento y regla de elegibilidad}
La tasa de Bangladesh pasa de 19.1\% el 1 de julio a 25.5\% el 2 de julio, y se mantiene entre 18.6\%
y 25.8\% hasta el 17. El sintético no muestra ese salto. La brecha media del 1 al 17 de julio es +4.8 pp, y
reaparece del 14 al 27 de agosto (+4.5 pp). Dos hechos:
\begin{itemize}
\item \textbf{Composición.} 891 usuarios tratados (8.7\% de la muestra) tienen su primera actividad
productiva el 29 de junio o después. En los donantes esa proporción va de 2.1\% a 4.5\%. Entran porque
la elegibilidad se define hasta el 17 de julio. Explican 440 de los 1,035 usuarios activos diarios
adicionales, y del 14 de agosto al 3 de septiembre son casi el doble de activos que el resto (28.2\% frente a
15.1\%).
\item \textbf{Nivel.} Sin esa cohorte, la tasa también sube: de 16–18\% a fines de junio a 18–24\% del
2 al 10 de julio.
\end{itemize}
El contrafactual depende de cómo se trate ese período (tabla \ref{tab:m2}). Ajustado solo con mayo y
junio, el sintético subestima a Bangladesh en 6.8 pp durante los 17 días previos al tratamiento: falla
la validación fuera de muestra más cercana que hay.

\begin{table}[h]
\small
\caption{Diagnóstico: control sintético según la ventana de ajuste}
\label{tab:m2}
\begin{tabularx}{\textwidth}{@{}L{3.6cm}rrrrrY@{}}
\toprule
 & ATT & Contra- & Efecto & Brecha & Brecha & \\
Ventana de ajuste & (pp) & factual & relativo & 1–17 jul & 14–27 ago & Donantes con peso \\
\midrule
1 may – 17 jul (paper) & −13.26 & 15.28 & −86.8\% & +4.82 & +4.46 & Pakistán, Ghana, Kenia \\
1 may – 30 jun & −11.32 & 13.34 & −84.8\% & +6.84 & +6.13 & Marruecos, Pakistán, Filipinas, Jordania, Indonesia \\
20 jun – 17 jul & −14.70 & 16.72 & −87.9\% & +3.11 & +4.08 & Ghana, Kenia \\
1 jul – 17 jul & −15.57 & 17.59 & −88.5\% & +2.47 & +3.53 & Kenia, Ghana \\
\bottomrule
\end{tabularx}
\par\smallskip{\footnotesize \emph{Notas:} mismo estimador, mismos 14 donantes y misma ventana de
tratamiento que el manuscrito. El SDID archivado da −17.95 pp. Fuente:
\path{07_referee_diagnostics.py}.\par}
\end{table}

El efecto relativo es estable (−85\% a −89\%); el efecto en puntos porcentuales no. Y la afirmación
de que la brecha «narrowed quickly» depende de la base: frente al nivel de julio, las semanas +1 y +2
siguen muy por debajo.
\pide
\begin{enumerate}
\item Definir la elegibilidad en una ventana anterior al período de estimación, de modo que no dependa
de julio. Exige decidir entre extraer datos anteriores a mayo o partir el período actual.
\item Averiguar qué ocurrió a inicios de julio y documentarlo con fuentes [CITATION NEEDED]. Verificar
con la fecha de creación de las cuentas si la cohorte es de cuentas nuevas.
\item Reportar la tabla \ref{tab:m2} en el texto principal como incertidumbre de especificación.
\item Matizar las frases sobre recuperación rápida en las secciones 6.2, 9 y 10.
\end{enumerate}

\coment{M3}{Apagón frente a disturbios}
El tratamiento es un paquete. La sección 8.4 lo reconoce, pero sus cuatro argumentos no discriminan:
el momento del quiebre y la recuperación son igual de compatibles con un toque de queda que empieza y
termina en las mismas fechas, y el rezago de commits, visto en niveles, es muy pequeño (M1).
\pide
\begin{enumerate}
\item Una cronología con fuentes de todo lo que cambió esos días, no solo la conectividad
[CITATION NEEDED].
\item Un análisis horario alrededor de los cambios de conectividad, que no coinciden entre sí: corte
total el 18 de julio hacia las 15:00 UTC, banda ancha fija el 23 hacia las 13:00 UTC, móvil el 28
hacia las 09:00 UTC (Belson 2024). GH Archive tiene resolución horaria.
\item Usar el período del 1 al 17 de julio. Si ya había movilizaciones sin apagón, la actividad
\emph{subió}; eso contradice que los disturbios por sí solos reduzcan la actividad, y hoy no se usa.
\item Buscar episodios comparables: apagones sin disturbios, como el de Rajasthan que estudian
Gunsilius y Van Dijcke (2023), y disturbios sin apagón.
\end{enumerate}

\coment{M4}{Inferencia}
\begin{itemize}
\item El valor $1/15$ es una restricción elegida. GH Archive cubre todos los países y el texto no
justifica por qué 14.
\item El test de rango no está calibrado. En el event study Bangladesh es la unidad más extrema de las
15 en 4 de las 11 semanas previas (−5, −3, −2 y −1) y en 5 de las 6 posteriores. Lo que distingue a la
semana 0 es la magnitud (13.09 frente a un máximo previo de 6.10), no el rango.
\item Bangladesh tiene el peor ajuste previo de las 15 unidades: RMSPE de 3.26 frente a un máximo de
2.87 entre los donantes (Kenia).
\item Un placebo en el tiempo (30 de junio a 5 de julio) alcanza un ratio de RMSPE de 3.79 frente al
4.19 real.
\item El patrón semanal difiere (tabla \ref{tab:m4}). El sintético cae sábado y domingo; Bangladesh
cae sobre todo el viernes. Eso genera brechas sistemáticas de +1.9 y +2.5 pp en sábado y domingo.
\end{itemize}

\begin{table}[h]
\small
\caption{Diagnóstico: tasa media por día de la semana, 1 de mayo a 30 de junio (\%, fechas UTC)}
\label{tab:m4}
\begin{tabularx}{\textwidth}{@{}Yrrrrrrr@{}}
\toprule
 & Lun & Mar & Mié & Jue & Vie & Sáb & Dom \\
\midrule
Bangladesh & 16.43 & 17.32 & 17.69 & 16.34 & 14.76 & 15.75 & 16.08 \\
Sintético & 17.29 & 17.64 & 17.56 & 16.81 & 15.71 & 13.89 & 13.57 \\
Brecha & −0.86 & −0.32 & 0.13 & −0.47 & −0.95 & 1.86 & 2.51 \\
\bottomrule
\end{tabularx}
\par\smallskip{\footnotesize Excluye los feriados de Bangladesh. Fuente:
\path{07_referee_diagnostics.py}.\par}
\end{table}

\pide
\begin{enumerate}
\item Ampliar el grupo de donantes con una regla fijada de antemano (por ejemplo, todos los países por
encima de un mínimo de usuarios elegibles) y exclusiones basadas en datos de conectividad de terceros.
\item Añadir inferencia que use las 78 observaciones previas, como la conformal de Chernozhukov,
Wüthrich y Zhu (2021), y discutir el ajuste imperfecto con Ferman y Pinto (2021).
\item Agregar a frecuencia semanal o quitar el efecto de día de la semana antes de ajustar.
\item Presentar el DiD como una diferencia de medias sobre los mismos datos, no como confirmación
independiente.
\end{enumerate}

\coment{M5}{Medición}
\begin{itemize}
\item \textbf{Días UTC y día parcial.} El 18 de julio tiene quince horas con conexión y una tasa de
10.41\%. Del 19 al 22 la tasa es 0.20\%, 0.18\%, 0.25\% y 0.22\%. El efecto en esos cuatro días es
−14.47 pp, −98.6\%. El −86.8\% del titular subestima el efecto de un día completo sin conexión.
\item \textbf{Ubicación.} Ese 0.2\% residual, unas 22 cuentas por día, acota cuántas cuentas
«de Bangladesh» operan desde fuera. Es una validación del autorreporte que el paper no usa.
\item \textbf{Selección sobre el outcome.} La muestra exige actividad dentro del período previo. Eso
induce reversión a la media: los donantes caen 2.15 pp entre el período previo y el apagón.
\item \textbf{Definiciones.} Ver R2. \textbf{Outcomes secundarios.} Ver R1.
\end{itemize}
\pide Reportar el efecto del 19 al 22 de julio junto al de la ventana completa; reportar la actividad
residual como cota del error de ubicación; declarar el huso horario; y discutir la reversión a la
media, que se resuelve con el punto 1 de M2.

\coment{M6}{RDiT}
El corte está en el 17 de julio, un día antes del tratamiento y del lado tratado, con 15
observaciones. El resultado depende de esas dos elecciones (tabla \ref{tab:m6}). El repositorio guarda
además una grilla anterior de 16 especificaciones, entre −12.4 y −23.7 pp, que el manuscrito no
reporta. Usa otra muestra (10,329 cuentas; 174 no están en la del control sintético), cuya actividad
residual durante el apagón es 0.5\% y no 0.2\%.

\begin{table}[h]
\small
\caption{Diagnóstico: estimado del RDiT (pp) según corte y ancho de banda $h$, con error estándar HAC}
\label{tab:m6}
\begin{tabularx}{\textwidth}{@{}Yrrrrr@{}}
\toprule
Corte & $h=4$ & $h=5$ & $h=7$ & $h=10$ & $h=14$ \\
\midrule
17 de julio (paper: $h=7$) & −4.45 (3.45) & −7.30 (4.36) & −9.28 (5.03) & −13.70 (6.05) & −16.80 (5.64) \\
18 de julio & −11.09 (2.72) & −12.98 (3.12) & −17.49 (4.48) & −18.06 (3.94) & −19.52 (3.47) \\
19 de julio & −9.00 (2.49) & −11.54 (2.91) & −16.59 (4.24) & −16.76 (3.15) & −17.60 (2.99) \\
\bottomrule
\end{tabularx}
\par\smallskip{\footnotesize Mismo estimador que \path{analysis/rdit/01_reproduce_rdit.py}. No es una
propuesta de nueva especificación. Fuente: \path{05_definitions_rdit.py}.\par}
\end{table}

\pide Eliminarlo, o rehacerlo con datos horarios y corte en el inicio del apagón. Si se queda, reportar
la grilla completa y justificar el corte. No elegir la especificación por su resultado.

\coment{M7}{Recuperación sectorial}
No tiene nivel de contribución principal.
\begin{itemize}
\item La ficha del clasificador dice que cerca de la mitad de lo que etiqueta con confianza no
pertenece a ningún sector, y el paper no lo valida en repositorios de Bangladesh.
\item Las celdas son pequeñas: de 213 a 1,528 usuarios, con tasas diarias de 0.06\% a 0.44\%.
\item No hay grupo de control ni contraste formal de diferencias entre sectores. «Finance recovered
fastest» compara índices sin error estándar de la diferencia.
\item Las líneas base son inestables: \emph{Accommodation/Food} tiene 254.1 y 127.3 en las dos
primeras semanas base; \emph{Finance}, 193.6 en la semana −1.
\item La ventana contiene los días de R5.
\item «Primera corrida de dos semanas $\geq 80$» es un estadístico ad hoc.
\end{itemize}
\pide Moverlo al apéndice como descripción, o reemplazarlo por heterogeneidad según rasgos del usuario
fijados antes del apagón (por ejemplo, pertenencia a organizaciones o patrón de actividad en días
laborables), que no necesita clasificador.

\coment{M8}{Pregunta económica y encuadre}
Faltan tres cantidades que el paper ya puede calcular: la pérdida acumulada en días-usuario activos
hasta que la serie vuelve a su nivel; si hubo recuperación por encima del contrafactual que compense
lo perdido; y la persistencia, para la que la muestra sectorial llega a noviembre. El párrafo de
política de la sección 9 (mirrors locales, redundancia internacional) no se apoya en ningún resultado.
La frase «GitHub timestamps reveal all three margins» promete más de lo que la sección 7 muestra.

\subsection{Comentarios menores}

\begin{enumerate}
\item Pakistán pesa 47.6\%. Belson (2024), ya citado, documenta una disrupción nacional allí el 31 de
julio, dentro de la semana +2; el texto solo menciona la del 14 al 17 de julio.
\item Ghana (1,373 cuentas) y Jordania (446) son donantes pequeños y ruidosos; considerar un tamaño
mínimo.
\item La sección 4.3 presenta 86,065 pares repositorio-SHA como «analysis target» sin decir que antes
del 18 de julio es una muestra de 17,000 de 134,479 (12.6\%). Eso aparece recién en la sección 7.
\item La marca de tiempo del \emph{committer} no es una verificación independiente: sale del mismo
reloj local.
\item Al quitar los diez usuarios con más rezago obtengo $N=2{,}740$ y 23.76\%; la tabla dice 2,739 y
23.77\%. Es un empate en el orden; conviene fijar el criterio.
\item «The two indices are 94.5 and 91.3» (sección 6.5): el orden es \emph{Information} y luego
\emph{Finance}, al revés de como se nombran en el párrafo.
\item «Returned to within its pre-shutdown range in the first week» (sección 9): con un rango previo
de −4.9 a +6.1 pp, ese criterio es poco exigente.
\item La fila «July 17 holiday» de la Tabla 1 no tiene fuente.
\item \path{paper/thesis.txt} y \path{paper/thesis.docx} tienen otro título y no incluyen la sección
sectorial; decidir cuál es la versión vigente y retirar las demás.
\end{enumerate}

\section{Bibliografía}

La verificación se hizo con búsqueda web (Crossref por DOI y lectura directa de las fuentes); no la
repetí de forma independiente, así que cada punto trae su enlace. Las 25 entradas existen. Ninguna
clave citada falta en \path{references.bib}.

\textbf{Correcciones a \path{references.bib}.}
\begin{itemize}
\item \cod{cloudflare2024}: el autor es David Belson; falta la URL
\url{https://blog.cloudflare.com/q2-2024-internet-disruption-summary/}.
\item \cod{isoc2024}: la URL apunta al listado general. La específica es
\url{https://pulse.internetsociety.org/en/shutdowns/internet-suspended-in-bangladesh-18-july-2024/}.
\item \cod{githubevents}: la página actual ya no contiene el tope de 20 commits; GitHub cambió los
payloads en 2025. Citar una versión fechada, por ejemplo
\url{https://docs.github.com/en/enterprise-server@3.14/rest/using-the-rest-api/github-event-types}.
\item \cod{schachenhofer2023}: faltan el volumen 98 y el número de artículo 104101.
\item \cod{aceto2018}: el apellido es Pescapé. \cod{kreitmeir2024}: quitar el comentario «VERIFY».
\item \cod{aceto2018} y \cod{schachenhofer2023} solo entran por \cod{\textbackslash nocite}: citarlas
o quitarlas.
\end{itemize}

\textbf{Afirmaciones que la fuente respalda solo en parte.}
\begin{itemize}
\item Sección 2.1, «Internet Society Pulse independently records»: esa página se apoya en Cloudflare
Radar e IODA. Mejor «also records».
\item Sección 2.1, «some residential fixed networks returned later»: Belson no dice «residential».
\item Tabla 1, fila del 16 y 17 de julio: el bloqueo selectivo lo respalda Internet Society, no
Belson.
\item Sección 3, sobre Kreitmeir y Raschky: usan indicadores binarios porque las acciones son raras;
no hablan de distribuciones asimétricas.
\item Sección 6.3, «consistent with Hausman and Rapson»: paráfrasis laxa; el artículo pide cautela,
no esa regla.
\end{itemize}

\textbf{Literatura que un árbitro esperaría} (existencia verificada por DOI).
\begin{itemize}
\item Chernozhukov, Wüthrich y Zhu (2021), \emph{JASA} 116(536): 1849–1864.
\url{https://doi.org/10.1080/01621459.2021.1920957}
\item Abadie, Diamond y Hainmueller (2015), \emph{AJPS} 59(2): 495–510. Origen del ratio RMSPE y los
placebos en el tiempo. \url{https://doi.org/10.1111/ajps.12116}
\item Ferman y Pinto (2021), \emph{Quantitative Economics} 12(4): 1197–1221.
\url{https://doi.org/10.3982/QE1596}
\item Ben-Michael, Feller y Rothstein (2021), \emph{JASA} 116(536): 1789–1803.
\url{https://doi.org/10.1080/01621459.2021.1929245}
\item Conley y Taber (2011), \emph{REStat} 93(1): 113–125. \url{https://doi.org/10.1162/REST_a_00049}
\item Zhang y Zhu (2011), \emph{AER} 101(4): 1601–1615. Bloqueo de acceso como shock a contribuciones
en línea. \url{https://doi.org/10.1257/aer.101.4.1601}
\item McDermott y Hansen (2021), NBER WP 29598. GitHub como medida de trabajo en tiempo real.
\url{https://doi.org/10.3386/w29598}
\end{itemize}

\section{Plan de trabajo}

Prioridad 1: antes de cualquier otra cosa. Prioridad 2: necesario para enviar. Prioridad 3: limpieza.

{\small\renewcommand{\arraystretch}{1.35}
\begin{longtable}{@{}L{0.9cm}L{5.6cm}L{4.7cm}L{2.0cm}c@{}}
\toprule
ID & Qué hacer & Dónde & ¿Análisis nuevo? & Prior. \\
\midrule
\endhead
R1 & Reconstruir los outcomes secundarios solo con elegibles; regenerar tablas y texto; reportar rangos de conteos &
  \path{country_day_outcomes_enriched.parquet}; \path{results/final/}; sección 6.4; introducción & Recalcular & 1 \\
R2 & Definir outcome primario, elegibilidad, bots y huso horario & Sección 4; Tabla 2 & No & 1 \\
R3 & Hacer que tablas y figuras corran en un clon limpio & \cod{.gitignore}; \path{make_all_tables.py}; \path{make_all_figures.py} & No & 1 \\
R5 & Documentar el 4 al 7 de agosto; corregir nota de la Figura 1 y sección 6.5 & Figura 1; secciones 2.1 y 6.5 & No & 1 \\
M1 & Mecanismo por fecha de autoría, cohorte desde el corte total, niveles & Sección 7; Figura 4; resumen; introducción; secciones 9 y 10 & Sí & 1 \\
M2 & Elegibilidad en ventana previa; causa del salto de julio; tabla de sensibilidad & Secciones 4.1, 6.1, 6.2 y 8.1 & Sí; acordar antes & 1 \\
M4 & Más donantes con regla previa; inferencia conformal; frecuencia semanal & Secciones 5.1, 5.4 y 6.1 & Sí; acordar antes & 2 \\
M3 & Cronología con fuentes; análisis horario; usar 1 al 17 de julio & Secciones 2.1 y 8.4 & Sí & 2 \\
M5 & Efecto del 19 al 22 de julio; cota de ubicación; reversión a la media & Secciones 6.1 y 8.2 & Ya calculado & 2 \\
M8 & Pérdida acumulada; reencuadre de contribuciones & Resumen; introducción; secciones 9 y 10 & Cálculo simple & 2 \\
R4 & Guardar conteos sectoriales y resumen \cod{pre8} & \path{results/sector_recovery/} & No & 2 \\
R6 & Subir el código aguas arriba & \path{analysis/} & No & 2 \\
M6 & Eliminar o rehacer el RDiT; reportar la grilla & Secciones 5.3 y 6.3 & Depende & 3 \\
M7 & Sectorial al apéndice o reemplazo & Secciones 4.4, 5.5 y 6.5 & Depende & 3 \\
Bib. & Correcciones y literatura & \path{references.bib}; secciones 2 y 3 & No & 3 \\
R7 & \cod{requirements.txt}; README; manifest & Raíz del repo & No & 3 \\
Men. & Comentarios menores 1 a 9 & Varios & No & 3 \\
\bottomrule
\end{longtable}
}

\appendix
\section{Scripts de auditoría}

Están en \path{audit_scripts/} y solo leen el repositorio. Cada uno se ejecuta con
\cod{python <script>.py}; la salida de esta auditoría está en \path{audit_scripts/output/}.

\begin{table}[h]
\small
\begin{tabularx}{\textwidth}{@{}L{5.3cm}Y@{}}
\toprule
Script & Qué verifica \\
\midrule
\path{01_rebuild_country_day.py} & Reconstruye el panel país-día del outcome principal y la regla de elegibilidad \\
\path{02_secondary_sample_check.py} & Muestra que los outcomes secundarios usan numerador de todos los usuarios (R1) \\
\path{03_secondary_sc_reproduce.py} & Reproduce las tablas secundarias archivadas y calcula la versión corregida (R1) \\
\path{04_mechanism_reproduce.py} & Reproduce las tablas del mecanismo; cohortes por fecha y hora de autoría; niveles (M1) \\
\path{05_definitions_rdit.py} & Contenido de \cod{productive\_events} (R2); muestra y sensibilidad del RDiT (M6) \\
\path{06_text_claims.py} & Cruza cifras del texto con los archivos de resultados \\
\path{07_referee_diagnostics.py} & Ventana de ajuste, composición diaria del ATT, cohorte de julio, día de la semana (M2, M4, M5) \\
\path{08_context_checks.py} & Feriados, fases del análisis previo, calibración del test de rango (R5, M4) \\
\path{09_compare_figures.py} & Compara figuras versionadas y regeneradas \\
\bottomrule
\end{tabularx}
\end{table}

\end{document}
